\section{Introduction}

University campus gyms are commonly imagined as open health spaces accessible to all students, yet their actual accessibility is far from equal.
In free-weight and strength training areas in particular, female students enter less frequently, stay for shorter durations, and exhibit more pronounced avoidance.
A recent survey of 357 female college students found that respondents reported moderate levels of gym intimidation on the \textcite{levinson2013seam} subscale, and that this intimidation had become a significant barrier to exercise participation; the same study also noted higher avoidance among lower-year students \parencite{rendall2024gym}.
Consistent with this finding, qualitative research on female gym users has shown that the intense gender segregation, masculine performance, sense of being watched, and spatial atmosphere of ``always being on display'' in the weight area can transform what should be a public training environment into a threshold that requires considerable courage to cross \parencite{turnock2021gyms}.

If campus gyms impose a form of physical spatial pressure, image-based social media platforms continuously generate a digital form of body pressure.
Existing research indicates that the link between social media and body image concerns does not operate primarily through ``screen time'' per se, but rather indirectly through appearance ideal internalization, appearance comparison, and appearance-related self-presentation \parencite{fardouly2016social}.
Research on young women further demonstrates that the relationship between social media use intensity and negative body image is mediated through the ``ideal internalization--social comparison'' chain \parencite{jung2022social}; experimental studies have also found that even fitspiration images, which ostensibly promote ``health and discipline,'' can immediately elevate body dissatisfaction and negative mood \parencite{tiggemann2015fitspiration}.

For the present study, what is more critical is that these two forms of pressure are not independent of each other.
Research on Chinese female college students has shown that media and peers influence negative body image through appearance comparison and thin-ideal internalization \parencite{shen2022media}; in a sample of young Chinese women, selfie posting and body comparison on social networks were significantly associated with body surveillance, weight-control exercise, and restrictive eating \parencite{yao2020selfie}.
In other words, many women are not simply ``unwilling to exercise''; rather, they are simultaneously navigating two sets of pressures:
on one hand, the free-weight area appears to be someone else's territory; on the other hand, digital media constantly reminds them that a woman's body should ideally remain slender and aesthetically pleasing, with some women also expressing concerns about appearing ``too muscular'' from strength training \parencite{peters2019barriers}.
Narrative research on Chinese women's fitness experiences also reminds us that fitness can be both a practice of agency and, simultaneously, a form of bodily discipline shaped by aesthetic norms \parencite{xiong2021fanshen}.

Building on this foundation, the present study focuses on resistance training participation among female college students in Chinese universities, attempting to answer a question with broader sociological significance:
When physical exercise spaces are experienced as male-dominated, and digital aesthetic standards continuously reinforce the norms of a ``correct female body,'' how do women understand, choose, and avoid certain forms of training?
The overarching goal of this study is to examine how the spatial environment of campus gyms and social media appearance norms jointly shape female college students' attitudes, preferences, and participation behaviors regarding resistance training.
Although the sample is drawn from a single Chinese university, the theoretical framework is intended to have broader relevance to Asian campus contexts where similar gendered spatial dynamics and digital beauty norms operate.
Around this goal, this paper poses three research questions:

\begin{enumerate}[label=RQ\arabic*:, leftmargin=*]
  \item How do the spatial environment of campus gyms and gym intimidation influence female college students' resistance training avoidance?
  \item To what extent does the internalization of social media appearance norms shape female students' exercise preferences?
  \item How do female students negotiate between ``exercising for health'' and ``fear of deviating from digital appearance norms''?
\end{enumerate}

This study carries dual significance.
Theoretically, it does not simply explain women's avoidance of campus gyms as a lack of personal motivation; instead, it places space, gaze, comparison, appearance internalization, and training choices within a single analytical framework, thereby addressing the fragmentation in existing research that examines ``body image only'' or ``exercise participation only.''
Practically, this issue directly concerns whether the allocation of campus physical education resources has genuinely achieved equitable accessibility:
if female students nominally have access to the gym but are effectively excluded from the free-weight area by spatial layout, body anxiety, and appearance concerns, then ``campus health promotion'' remains incomplete.

The structure of this paper is as follows: the next section reviews the relevant literature on gendered gym spaces, resistance training barriers, body objectification, and social media appearance norms; the subsequent section describes the methods and data limitations of this study; the results section presents the questionnaire aggregate data and interview themes; the discussion section is organized around four themes---spatial pressure, beginner barrier, appearance influence, and strategic negotiation---with responses to the three research questions embedded within each subsection.
